% Section 2.4, from lectures/L02/appendix/b_exercises.md. The solutions are in the repository, in
% lectures/L02/exercises.

\section{Exercises}
\label{c2:sec:exercises}\label{c2:app:b}

\begin{aside}
\textbf{How to work these.} All four sets extend one program: the linear regression model you
wrote in Chapter~\ref{c1:ch:linreg}'s Set~2. Carry it forward into this lecture's exercises
directory, \file{lectures/L02/exercises} in the companion repository, and extend it set by set,
rebuilding with \code{make} after each exercise.

\textbf{How to check your work.} Set~4 runs the lecture's test suite, already in
\file{lectures/L02/exercises/test}. It is cumulative: it covers everything from
Chapter~\ref{c1:ch:linreg} as well as what this chapter adds, so it supersedes the earlier suite
entirely. It tests the directory above it by default; \code{make ML_DIR=<path>} points it at a tree
of your own elsewhere.

The solutions are published in the companion repository, in \file{lectures/L02/exercises}, with a
README that walks through them, including what the adaptive learning rate buys on training data its
initial rate is too cautious for. Try each exercise yourself before looking at them.
\end{aside}

You'll extend the \code{ml::lin_reg::Fixed} class from Chapter~\ref{c1:ch:linreg}, first with a
randomized training order, then with precision calculation. Finally you'll add a second model,
\code{ml::lin_reg::Adaptive}, which picks its own learning rate while it trains.

% ------------------------------------------------------------------------------------------
\exerciseset{Randomizing the Training Order}
\label{c2:set:1}

\codeexercise{MatrixU32 alias}
\label{c2:ex:1}
Add the alias \code{MatrixU32} to the header file \file{ml/types.hpp} as a substitute for
\code{std::vector<std::uint32_t>}, in the namespace \code{ml}.

Include \code{<cstdint>} in \file{ml/types.hpp}.

\codeexercise{Random number generator}
\label{c2:ex:2}
Add a function named \code{initRandGen()} in an anonymous namespace in
\file{source/ml/lin_reg/fixed.cpp}. The function should initialize the random number generator
exactly once when called:
\begin{itemize}
  \item Use a static local variable named \code{initialized} (of type \code{bool}, initialized to
    \code{false}) to track whether the generator has already been initialized.
  \item Initialize the generator with \code{std::srand(std::time(nullptr))}.
  \item Set \code{initialized} to \code{true} after initialization.
  \item Include \code{<cstdlib>} and \code{<ctime>} in \file{fixed.cpp}.
  \item Should be marked \code{noexcept}.
\end{itemize}

\codeexercise{Training order}
\label{c2:ex:3}
Add a private member variable named \code{myTrainOrder} to \code{Fixed}:
\begin{itemize}
  \item Should hold the indices of the training sets as unsigned integers (\code{ml::MatrixU32}).
  \item Mustn't be marked \code{const}, since its contents are rearranged before every epoch.
\end{itemize}
The member variable \code{mySetCount} from Chapter~\ref{c1:ch:linreg} can now be removed:
\code{myTrainOrder} holds one index per training set, so \code{myTrainOrder.size()} is the number of
training sets.

Update the constructor in \file{source/ml/lin_reg/fixed.cpp}:
\begin{itemize}
  \item Compute the number of complete training sets as in Chapter~\ref{c1:ch:linreg}, i.e.\
    \code{std::min(trainIn.size(), trainOut.size())}, but keep it in a local variable rather than a
    member.
  \item Print an error message and call \code{std::terminate()} if the set count is 0, as in
    Chapter~\ref{c1:ch:linreg}.
  \item Call \code{initRandGen()}, so the random number generator is ready before the first
    shuffle. Calling it here rather than in \code{train()} means it's initialized exactly once per
    model, no matter how often training is started.
  \item Resize \code{myTrainOrder} to the set count, then fill it with the indices
    \code{0, 1, 2 ... N-1}.
\end{itemize}

\codeexercise{Shuffling}
\label{c2:ex:4}
Add a private method named \code{shuffle()} to \code{Fixed}:
\begin{itemize}
  \item Should shuffle the contents of \code{myTrainOrder} in random order.
  \item For each index \code{i}, pick a random index \code{r} and swap \code{myTrainOrder[i]} and
    \code{myTrainOrder[r]}.
  \item Should be marked \code{noexcept}.
\end{itemize}

\codeexercise{Updating the training method}
\label{c2:ex:5}
Update the method \code{train()} in \file{source/ml/lin_reg/fixed.cpp}:
\begin{itemize}
  \item Call \code{shuffle()} at the start of every epoch.
  \item Iterate through the training sets in the order given by \code{myTrainOrder} instead of
    sequentially by index, e.g.\ with a range-based for loop over \code{myTrainOrder}. Note that the
    loop variable is now the index into the training data, not the counter itself.
\end{itemize}

\codeexercise{Compiling and testing}
\label{c2:ex:6}
Compile and test-run the program. The training order should now be random every epoch, but the
predicted result should still match the training data once training is complete.

% ------------------------------------------------------------------------------------------
\exerciseset{Precision Calculation}
\label{c2:set:2}

\codeexercise{Precision calculation}
\label{c2:ex:7}
Add a private method named \code{precision()} to the \code{Fixed} class. The method should compute
the model's precision given the training data:
\begin{itemize}
  \item Compute the mean absolute error (MAE) across all training sets:
  \begin{itemize}
    \item For each training set, compute the absolute error: \code{abs(output - predict(input))}.
    \item Sum all absolute errors and divide by the number of training sets.
  \end{itemize}
  \item Return \code{1.0 - MAE} as a floating-point number.
  \item Should be marked \code{[[nodiscard]]}, \code{const}, and \code{noexcept}.
\end{itemize}
Include \code{<cmath>} in \file{fixed.cpp} for \code{std::abs}.

\codeexercise{Updating the training method}
\label{c2:ex:8}
Update the method \code{train()} in \file{source/ml/lin_reg/fixed.cpp} so training stops early once
the precision exceeds a given threshold:
\begin{itemize}
  \item Add an argument named \code{precisionThreshold} (a floating-point number specifying the
    precision threshold) as the \textbf{last} parameter, so the full signature reads
    \code{train(epochCount, learningRate = 0.01, precisionThreshold = 0.999999)}. Default value:
    \code{0.999999} (99.9999\%). Note that the argument shouldn't be named \code{precision}: inside
    \code{train()}, a parameter of that name hides the method \code{precision()}, so calling it no
    longer compiles (unless \code{this->precision()} is used).
  \item Return \code{false} if \code{precisionThreshold} falls outside the range
    \code{(0.0, 1.0)}:
  \begin{itemize}
    \item A threshold of \code{1.0} or more can never be reached, since the precision is
      \code{1.0 - MAE} and the mean absolute error can't be negative. Training would always run the
      full epoch count.
    \item A threshold of \code{0.0} or less accepts a model whose mean absolute error is \code{1.0}
      or worse, which defeats the purpose of the check.
  \end{itemize}
  \item As in Chapter~\ref{c1:ch:linreg}, \code{train()} reports invalid arguments through its
    return value. Only the constructor calls \code{std::terminate()}, because it has no way to
    return a failure code to the caller.
  \item Call \code{precision()} every tenth epoch, starting with the first; computing it scans every
    training set, so checking less often than every epoch keeps that overhead down over thousands
    of epochs. Use a named constant for the interval rather than a bare \code{10}.
  \item Stop training and return \code{true} as soon as the precision reaches or exceeds
    \code{precisionThreshold}, printing the achieved precision and the number of epochs it took.
  \item Return \code{true} once the epochs run out, if the threshold was never reached.
\end{itemize}

\codeexercise{Compiling and testing}
\label{c2:ex:9}
Compile and test-run the program. The model should now stop training as soon as the precision
reaches the threshold, long before the 1000 epochs \file{main.cpp} asks for:

\begin{console}
Target precision 1.00 reached after 11 epochs!
Input: 0, prediction: 2
Input: 1, prediction: 5
Input: 2, prediction: 8
Input: 3, prediction: 11
Input: 4, prediction: 14
\end{console}

\noindent The training order is random, so the epoch count can vary between runs. For this
training data it settles at 11: the precision is only checked every tenth epoch, so the check that
first passes is the one after the eleventh epoch has run. Report the number of epochs completed, not
the loop index, so the first epoch reads as 1 rather than 0.

% ------------------------------------------------------------------------------------------
\exerciseset{An Adaptive Learning Rate}
\label{c2:set:3}
The learning rate is the one training parameter with no sensible default. Too low and the model
crawls, too high and it never settles, and where the line between the two falls depends on the
training data rather than on the algorithm. Instead of asking the caller to guess, you'll now write
a second model that starts at a cautious rate and revises it while training, using the precision it
already computes every tenth epoch.

The new class is called \code{Adaptive}, and it's the one the \code{Fixed} name has been pointing at
since Exercise~\ref{c1:ex:7}.

\codeexercise{Shared utilities}
\label{c2:ex:10}
\code{initRandGen()} currently sits in an anonymous namespace in \file{fixed.cpp}, which makes it
private to that file. The new model needs the same one-time seeding, and copying the function across
would give each file its own \code{initialized} flag: the second copy would call \code{std::srand()}
again when the second model is constructed, and a re-seed within the same second restarts the very
sequence that was already running.

Move the function into a pair of shared files instead:
\begin{itemize}
  \item Create \file{include/ml/utils.hpp}. Declare \code{initRandGen()} there, in the namespace
    \code{ml}, taking no arguments, returning nothing, and marked \code{noexcept}.
  \item Create \file{source/ml/utils.cpp}. Move the definition there unchanged, along with the
    \code{<cstdlib>} and \code{<ctime>} includes it needs.
  \item Remove the function and the \code{<ctime>} include from \file{fixed.cpp}, and include
    \code{"ml/utils.hpp"} instead. Keep \code{<cstdlib>}: \code{shuffle()} still calls
    \code{std::rand()}. The call in the constructor stays exactly as it was: \code{Fixed} lives in
    \code{ml::lin_reg}, so the unqualified name still finds \code{ml::initRandGen()} in the
    enclosing namespace.
  \item Add \file{source/ml/utils.cpp} to \code{SRC_FILES} in the makefile.
\end{itemize}
Rebuild and re-run before going further. Nothing about the program's behaviour should have changed;
this step only moves code.

\codeexercise{The Adaptive class}
\label{c2:ex:11}
Add the header \file{include/ml/lin_reg/adaptive.hpp} and the source file
\file{source/ml/lin_reg/adaptive.cpp}, and add the source file to \code{SRC_FILES} in the makefile.

In the namespace \code{ml::lin_reg}, implement a class named \code{Adaptive} that inherits
\code{Interface} via public inheritance and is marked \code{final}. Everything except
\code{train()} is the same as in \code{Fixed}:
\begin{itemize}
  \item The constructor takes \code{trainIn} and \code{trainOut} and is marked \code{explicit} and
    \code{noexcept}. It checks the set count, terminates when it's 0, fills \code{myTrainOrder}
    with \code{0, 1, 2 ... N-1}, and calls \code{initRandGen()}.
  \item \code{~Adaptive()} is \code{default}, \code{noexcept}, and \code{override}.
  \item \code{predict()} returns $y = kx + m$, and is marked \code{[[nodiscard]]}, \code{const},
    \code{noexcept}, and \code{override}.
  \item The private methods \code{optimize()}, \code{shuffle()}, and \code{precision()} are the same
    as in \code{Fixed}.
  \item The private member variables are the same: \code{myTrainOrder}, \code{myTrainIn},
    \code{myTrainOut}, \code{myBias}, and \code{myWeight}.
  \item The default constructor, copy and move constructors, and the corresponding operators are
    deleted.
\end{itemize}
\code{train()} is where the two models part ways:

\begin{cppcode}
bool train(std::size_t epochCount, double precisionThreshold = 0.999999) noexcept;
\end{cppcode}

\noindent There's no \code{learningRate} argument, because the model no longer takes one. Copying
the shared parts across rather than sharing them with \code{Fixed} is deliberate here: two small
independent models make the one difference between them easy to see.

\codeexercise{The learning rate rule}
\label{c2:ex:12}
Add a function named \code{updateLearningRate()} in an anonymous namespace in
\file{adaptive.cpp}:

\begin{cppcode}
void updateLearningRate(double& learningRate, double& prevPrecision,
                        double currentPrecision) noexcept;
\end{cppcode}

\noindent It takes the current learning rate and the previous precision by reference, since it
updates both, and the precision just measured by value. Give each of the four values it needs a
named constant: a maximum learning rate of \code{0.25}, a minimum of \code{1e-6}, a smallest
expected improvement of \code{0.1}, and a step of \code{0.05}.

The rule itself works on the difference between the two precisions, i.e.\ how much the model
improved since the last evaluation:
\begin{itemize}
  \item \textbf{Improved by at least the expected minimum:} leave the learning rate alone. It's
    working.
  \item \textbf{Improved, but by less than that:} raise the learning rate by the step. Progress has
    stalled, so take bigger steps.
  \item \textbf{Didn't improve at all:} lower the learning rate by the step. The model is
    overshooting, so take smaller steps.
\end{itemize}
Neither branch should enforce the limits itself. Once the rate has been adjusted, keep it inside
them with a single call to \code{std::clamp()} (from \code{<algorithm>}, already included for
\code{std::min()}):
\code{learningRate = std::clamp(learningRate, minLearningRate, maxLearningRate)}. It returns the
nearer bound whenever the value falls outside the range, so one call covers both ends and each
branch is left saying only what it's about.

Store the precision just measured in \code{prevPrecision} before returning, so the next call has
something to compare against.

That one clamp is what keeps the rule safe. Its maximum stays below the rate at which this kind of
training data starts to oscillate instead of converging, and its minimum keeps the rate positive, so
a model that keeps missing can still creep towards the line rather than stopping dead.

\codeexercise{Training}
\label{c2:ex:13}
Implement \code{train()} in \file{adaptive.cpp}:
\begin{itemize}
  \item Give the evaluation interval (\code{10}) and the initial learning rate (\code{0.1}) named
    constants.
  \item Declare the learning rate and the previous precision as local variables, initialized to the
    initial learning rate and \code{0.0}. Being local means every call starts from the initial rate
    again, rather than inheriting whatever the previous call ended on.
  \item Return \code{false} if \code{epochCount} is 0, or if \code{precisionThreshold} falls outside
    the range \code{(0.0, 1.0)}, exactly as \code{Fixed} does. There's no learning rate left to
    validate.
  \item For each epoch: call \code{shuffle()}, then run \code{optimize()} over all training sets in
    the order given by \code{myTrainOrder}, passing the current learning rate.
  \item Every tenth epoch, except the first, compute the precision once and keep it in a local
    variable:
  \begin{itemize}
    \item Stop training and return \code{true} if it reaches or exceeds \code{precisionThreshold},
      printing the achieved precision and the number of epochs it took, as in
      Exercise~\ref{c2:ex:8}.
    \item Otherwise, pass it to \code{updateLearningRate()}.
  \end{itemize}
  \item Return \code{true} once the epochs run out.
\end{itemize}
Note that the precision is computed once per evaluation and then used twice. Calling
\code{precision()} a second time inside \code{updateLearningRate()} would scan every training set
again for a value that's already in hand.

\codeexercise{Compiling and running}
\label{c2:ex:14}
Update \file{main.cpp} to train an \code{ml::lin_reg::Adaptive} model rather than a \code{Fixed}
one, and drop the learning rate from the \code{train()} call. An epoch budget of 100 is plenty now,
since the model stops as soon as it reaches the default precision threshold.

Compile and run. The predictions should match the training data, without a learning rate having
been chosen anywhere in the program:

\begin{console}
Input: 0, prediction: 2
Input: 1, prediction: 5
Input: 2, prediction: 8
Input: 3, prediction: 11
Input: 4, prediction: 14
\end{console}

\noindent Above them sits the early-stop line your own \code{train()} prints. The epoch count in it
varies with the shuffled order.

% ------------------------------------------------------------------------------------------
\exerciseset{Running the Tests}
\label{c2:set:4}

\codeexercise{Running the tests}
\label{c2:ex:15}
An updated test suite is available in \file{lectures/L02/exercises/test}. It's cumulative: it
covers everything from Chapter~\ref{c1:ch:linreg} as well as the randomized training order, the
precision threshold, and the adaptive learning rate added here, so it supersedes the earlier suite
entirely.

Once you've carried your code forward into this lecture's exercises directory, build and run it as
before:

\begin{shell}
make -C test
\end{shell}

\noindent All 38 test cases should pass: 20 for \code{Fixed} and 18 for \code{Adaptive}. Note that
\code{train()} prints its progress, so the output is interleaved with early-stop reports. Look for
the summary on the last line.

The test suite's README, \file{lectures/L02/exercises/test/README.md}, has more information,
including why the tests that check exact values train on a single training set now that the order
is randomized, and how a learning rate the model keeps to itself can be tested at all.
